% !TeX root = ../main.tex
\section{Long-strip detector}
\label{sec:long-strips}
The outer tracker uses complete stereo pairs on a continuous cooled structural
sandwich. Two tilted barrel layers and six petal discs at each end share a
one-dimensional sensor contract. The sandwich between the sensors provides
their local mount, thermal path and stave or petal structure; no additional
unmodeled rail is credited with carrying the load.

\subsection{Stereo sensors and module construction}
Each face has a \qty{96}{\milli\metre}-square active area in
\qty{0.3}{\milli\metre} silicon, with \qty{0.5}{\milli\metre} guard and
1200 strips at \qty{80}{\micro\metre} pitch. The physical in-plane axes rotate
by $\pm\qty{20}{\milli\radian}$, giving a \qty{40}{\milli\radian} relative
stereo angle. Each face returns one measured coordinate. Combining a pair must
transport both measurements to a common reference plane using the track slope;
the finite face separation must not be hidden in a fictitious two-coordinate
segmentation.

The barrel core is \qty{5}{\milli\metre} carbon foam with
\qty{0.10}{\milli\metre} adhesive and \qty{0.20}{\milli\metre} CFRP skins on
both sides. Each sensor couples through \qty{0.20}{\milli\metre} module glue,
a full-area \qty{0.10}{\milli\metre} graphite spreader and a
\qty{42}{\milli\metre}-square pickup. The pickup is
\qty{0.10}{\milli\metre} high on a low row and
\qty{1.60}{\milli\metre} on a raised row. Complete pairs are lifted outward
on both faces, giving midplane separations of \qty{6.7}{\milli\metre} or
\qty{9.7}{\milli\metre}. This replaces the thinner analytical gap with the
actual occupied sandwich. Electrical hybrids and insulated buses occupy the
opposite edge from the mounting web.

\subsection{Barrel placement and bearing connection}
Stave anchors at $r=840$ and \qty{1060}{\milli\metre} carry 86 and 108 staves,
respectively. Each stave has 36 complete pairs, giving \num{6984} pairs and
\num{13968} independent sensor faces. The common azimuthal tilt is
$+\qty{12}{\degree}$. The selected \qty{32}{\milli\metre} transverse reserve
and \qty{70.765}{\milli\metre} row pitch retain projected overlap for finite
stereo separation and curved incidence. They increase silicon and channel
inventory substantially relative to the original sparse layout; this is an
explicit coverage/material tradeoff. Full finite-plane screens, rather than
body projections alone, determine whether both faces of the same pair are hit.

\begin{figure}[tbp]
\centering\includegraphics[width=\linewidth]{long-strip-barrel}
\caption{True long-strip barrel section at global $z=166.50$ mm, inside a
bearing station. Left: both layers. Right: the innermost stave, bearing ring
and local mounting connection at equal transverse scales. All shapes are
intersections of the retained solids, including the chord of the potted
transverse pin. Purple is sensitive silicon, teal the bearing and mount,
grey structural and thermal components, and gold local electrical services.}
\label{fig:long-barrel}
\end{figure}

The stave is \qty{124}{\milli\metre} wide and extends to
$|z|=\qty{1320}{\milli\metre}$. The free shingle edge at local
$u=+\qty{55}{\milli\metre}$ carries its mounting lands; local buses and
hybrids lie at $u=-\qty{55}{\milli\metre}$. Two potted epoxy/titanium
fasteners at every station connect through a keyed CFRP web and tangent shoe to
an inner CFRP bearing ring. Rings occupy $r-45$ to
$r-42\,\unit{\milli\metre}$ and are \qty{6}{\milli\metre} wide axially.
The centre station locates $z$; the remaining stations propose longitudinal
sliding with $\pm\qty{0.3}{\milli\metre}$ travel. This is an assembly provision
for thermal motion and placement allowance, not a detailed machined slot or a
qualified clamp/preload specification.

\subsection{Anchors and sag limits}
The structural screen includes the attached core, skins, sensors, thermal
parts, electronics, tubes, buses and local mounting load. Grounded bearing rings
are excluded from the stave payload; the worst stave carries about
\qty{1.8}{\kilogram}. A released-span bending-and-core-shear calculation
uses skin modulus \qty{70}{\giga\pascal}, core shear modulus
\qty{5}{\mega\pascal} and twice gravity. These are conservative screen
hypotheses within the tested range, not measured material lower bounds.

The allocated normal displacement is \qty{50}{\micro\metre}, divided into
\qty{40}{\micro\metre} gravity deformation and a
\qty{10}{\micro\metre} joint allowance. Fourteen stations give
\qty{42.25}{\micro\metre} gravity deformation and fail the allocated gravity
budget; fifteen give \qty{35.49}{\micro\metre} and are the minimum under this
screen. \textbf{17 stations per stave}, at
\qty{163.75}{\milli\metre} spacing from $z=-1310$ to
$+1310\,\unit{\milli\metre}$, including the centre, are realised.

\begin{figure}[tbp]
\centering\includegraphics[width=\linewidth]{long-strip-anchors}
\caption{Selected long-strip bearing schedule and retained normal-load screen.
The released-span estimate includes bending, foam shear and overhang allowance.
The separate continuous Timoshenko solution is shown at 17 stations. The curve
uses the stated payload and elastic hypotheses; it gives no qualification of
the actual joints, rings, thermal bow or torsion.}
\label{fig:long-anchors}
\end{figure}

At 17 stations the bounded result is \qty{25.95}{\micro\metre}, comprising
\qty{5.06}{\micro\metre} bending, \qty{20.67}{\micro\metre} core shear and
a small overhang allowance. An independent continuous Timoshenko model gives
\qty{23.46}{\micro\metre}, with the two retained mesh spacings agreeing within
\qty{0.1}{\micro\metre}. Adding the joint allocation leaves margin within the
\qty{50}{\micro\metre} fixture goal. Core shear dominates at this spacing;
crediting only skin bending would understate the required support count.
Actual bond and fastener compliance, ring/global-frame deformation, local
indentation, torsion, vibration and thermal bow require independent qualification.

\subsection{Barrel cooling and cable routes}
Each half-stave has two U-circuits serving nine complete pairs each. Four
straight titanium legs at $u=-30,-10,+10,+30\,\unit{\milli\metre}$ use
\qty{2.5}{\milli\metre} outside diameter and \qty{0.14}{\milli\metre} wall.
Their separated returns have \qty{10}{\milli\metre} centre radius near
$|z|=\qty{20}{\milli\metre}$. Pipe envelopes are explicitly removed from the
foam. Supply and return emerge at the outer signed stave end. Both sensor faces
connect through insulated local buses and end boards to two harnesses per
half-stave; harness assignment counts complete pairs once.

The conductor-sized cable fixture allocates a \qty{6}{\milli\metre} LV/HV
bundle and a \qty{3.6}{\milli\metre} bundle containing 48 coated fibres.
The power envelope contains \qty{6.14}{\milli\metre\squared} copper and
\qty{5.24}{\milli\metre\squared} insulation per unit length. These conductor
and volume hypotheses are distinct from the larger gross cable comparator used
for short strips. The \qty{11}{\volt}, \qty{88.8}{\watt},
\qty{4.1}{\metre} comparator gives loop/return drops of
\qty{0.397}{\volt}/\qty{0.199}{\volt}; increasing resistivity by a factor 1.4
fails the return criterion. Dielectric, terminations, leakage and operational
power still need a component-level specification.

The barrel collector occupies $r=800$--$1169\,\unit{\milli\metre}$,
$|z|=1325$--$1435\,\unit{\milli\metre}$. Its downstream trunk uses the fixed
$r=1144$--$1169\,\unit{\milli\metre}$ corridor to
$|z|=\qty{3500}{\milli\metre}$. Cable and pipe turns reserve
\qty{50}{\milli\metre} and \qty{25}{\milli\metre} centre radii.
These are clearance reservations, without qualified elbows, weld access,
manifold balancing or connector installation. The barrel-only conductor-sized
packing cases pass; the larger \qty{13.4}{\milli\metre} gross power comparator
does not. Endcap demand must subsequently be included in the same allocation.

\subsection{Endcap petals, mounting and stereo placement}
The six discs per end use the long-strip datums in \cref{tab:strip-discs}.
The nominal active annulus is $r=786.46$--$1108.33\,\unit{\milli\metre}$.
Four rings contain $72,72,84,84$ complete pairs, or 312 pairs per disc and 26
per petal. Twelve petals per disc give \num{3744} pairs and
\num{7488} faces across both ends. Each pair keeps the tangential/radial
reference axes and physical $\pm\qty{20}{\milli\radian}$ stereo rotations.

The petal uses the same \qty{5}{\milli\metre} foam core, with
\qty{0.10}{\milli\metre} glue and \qty{0.30}{\milli\metre} skins.
Four full-height integral CFRP webs replace corresponding foam volumes;
they are part of the carrier, not an added backing plate. Four outward
elevations give sensor-face separations of $6.9,9.9,12.9,15.9\,\unit{\milli\metre}$.
The carrier spans $r=784$--$1130\,\unit{\milli\metre}$. Inner and outer
bearing rings span $783.5$--$785.5$ and $1130$--$1140\,\unit{\milli\metre}$.
The inner interface leaves only \qty{0.5}{\milli\metre} to the short-strip
service boundary and needs a detailed clearance/tolerance review.

\begin{figure}[tbp]
\centering\includegraphics[width=\linewidth]{long-strip-endcap}
\caption{Long-strip endcap. Left: one sensor face of each pair on the first
positive disc. Right: one continuous structural petal, with its four integral
webs, twelve straight legs of six cooling circuits and three mounting contacts.
All depths are projected. Blue marks tubes, grey integral webs and red the
mounting contacts. Both sensors of each pair have distinct physical
stereo frames; only one face is drawn to expose the carrier and routing.}
\label{fig:long-endcap}
\end{figure}

Potted lands and keyed shoes carry load to three contacts. One outer contact
fixes in-plane position, the other slides tangentially, and the inner key slides
radially. The $\pm\qty{0.3}{\milli\metre}$ motion provision avoids redundant
rigid thermal constraints. The proposed contacts, narrow inner key and grounded
rings require plate/joint/global-frame FEA and cold-cycle testing.

The worst complete petal payload is about \qty{1.5}{\kilogram}. At assumed
$E=\qty{70}{\giga\pascal}$ and $G=\qty{5}{\mega\pascal}$, the
\num{0.25}$g$ axial fixture gives \qty{43.85}{\micro\metre} including the
\qty{10}{\micro\metre} joint allocation. A horizontal handling load of $1g$
gives \qty{145.41}{\micro\metre} and fails even the tested stiffer combination.
A temporary handling frame is therefore recommended. This beam proxy assumes
distributed rim reactions; three discrete contacts, plate twisting, local
indentation and in-plane gravity cannot be qualified by that result.

\subsection{Petal cooling, dielectric and cumulative services}
Six U-loops per petal distribute twelve straight legs from local radius
$805$ to \qty{1105}{\milli\metre}. Tube size and
\qty{10}{\milli\metre} returns match the barrel. Nearest-leg assignment gives
$4,4,6,4,4,4$ complete pairs per circuit, counting both faces as one heat-load
unit. The integral webs stop before $r=\qty{1114}{\milli\metre}$; outer
collection buses start at \qty{1115}{\milli\metre}. Core holes and bus cutouts
are disjoint, conserving the remaining structural material.

The final \qty{0.05}{\milli\metre} copper foil has
\qty{0.15}{\milli\metre} polyimide around every exterior face. The spine and
outer collector share their intended conductor junction; dielectric is
partitioned around that joint so it is neither a bare-carbon contact nor a
double-subtracted foam region. This geometric correction does not qualify
breakdown voltage, radiation endurance or termination practice.

Three outer harnesses per petal serve groups of $12,12,2$ pairs. They reach a
$r=1100$--$1169\,\unit{\milli\metre}$ fan, at local outward offsets
$20$--$130\,\unit{\milli\metre}$, before the longitudinal trunk. The cumulative
endcap service inventory includes the pinned barrel loads and every upstream
disc; it replaces the standalone barrel corridor. The final reference demand
is \qty{3469.83}{\milli\metre\squared} per sector against
\qty{5449.88}{\milli\metre\squared} capacity. 

The farthest pair centre is \qty{30}{\milli\metre} from a pipe leg.
At \qty{3.7}{\watt} per face, normal foam-slab lower bounds are small, but a
pessimistic single lateral path gives temperature rises of
\qtyrange{26.43}{105.71}{\kelvin} over the tested conductivity range and fails
all temperature fixtures. The lower bound omits spreading/contact resistance;
the lateral case omits beneficial parallel paths. Neither is a predicted sensor
temperature. Anisotropic thermal analysis, qualified graphite-to-pipe bridges
or a revised pipe distribution must establish the operating margin. Enthalpy
budgets also omit pressure drop, two-phase stability and dry-out.